\section{Conclusão}

Este artigo apresentou um pipeline SNN multimodal offline com consciência energética para a predição da qualidade da água em um contexto amazônico de monitoramento na borda. A abordagem proposta demonstrou viabilidade prática do treinamento à implantação, incluindo interpretação geoespacial e benchmarking multiplataforma de latência/energia. Os resultados indicam que a inferência neuromórfica é implantável sob restrições reais de hardware, com compromissos mensuráveis e específicos por plataforma entre vazão e energia por inferência. Embora nem sempre seja a melhor em MAE absoluto, a SNN permanece uma opção robusta quando o processamento temporal orientado a eventos e a compatibilidade com a borda são requisitos de primeira ordem.

As ameaças à validade devem ser explicitamente consideradas na interpretação destes achados. O principal risco à validade interna é o tamanho limitado do conjunto de dados original de campo, que pode aumentar a sensibilidade à partição e tornar pequenas diferenças de métricas instáveis entre partições; isso foi mitigado por avaliação com sementes repetidas, aumento de dados conservador e relato por alvo, mas não totalmente eliminado. A validade externa também é limitada pela diversidade hidrológica entre sub-regiões amazônicas, de modo que afirmações regionais mais amplas requerem campanhas em múltiplos sítios e estações, com protocolos de aquisição harmonizados. Trabalhos futuros incluem aquisição de campo mais ampla ao longo de regimes sazonais, validação externa em sistemas de água doce \rev{e lagos} adicionais, e co-otimização mais profunda dos parâmetros de codificação e dos modos de execução do hardware para melhor equilíbrio energia-desempenho. \rev{As extensões metodológicas planejadas incluem validação cruzada k-fold completa no conjunto de dados ampliado, uma comparação sistemática de estratégias de aumento de dados além da perturbação gaussiana, a medição do consumo de energia de Random Forest, LSTM e MLP sob o mesmo protocolo de repouso/carga usado para a SNN em dispositivos idênticos, e modalidades de entrada enriquecidas para tratar a limitação do oxigênio dissolvido identificada na análise por alvo.} De forma geral, este estudo demonstra que modelos neuromórficos podem apoiar o monitoramento ambiental prático em regiões remotas, viabilizando inteligência de borda com consciência energética para a predição da qualidade da água na Amazônia.
